\chapter{Происхождение shift-симметрии Mathai--Quillen-пары}
\label{ch:v10-h15-r2-mathai-quillen-shift-symmetry-parent}

Предыдущий аудит свёл происхождение нечётной Thom-пары к одному вопросу:
имеет ли поле $\Sigma$ восьмимерную аддитивную gauge orbit
\begin{equation}
 \Sigma\longmapsto\Sigma+\varepsilon,
 \qquad \varepsilon\in\mathbb R^8,
 \qquad T_{\rm req}=I_8?
\end{equation}
Именно ghost этой symmetry мог бы быть полем $\psi_\Sigma$.

После условного устранения auxiliary-поля квадратичный оператор физической
$\Sigma$ равен
\begin{equation}
 G_Y=\operatorname{diag}(40,40,0,48,48,0,40,40),
 \qquad \rank G_Y=6,
 \qquad \dim\ker G_Y=2.
\end{equation}
Для линейной translation symmetry требуется $G_YT=0$. Но
\begin{equation}
 \rank(G_YT_{\rm req})=6,
\end{equation}
поэтому полная rank-$8$ orbit разрушает шесть квадратичных направлений.
Максимальный flat translation subspace quadratic parent имеет rank $2$ и
порождается двумя компонентами с $|6Y|=7$. Это всё ещё оставляет дефицит
$8-2=6$.

Обычная калибровочная симметрия не восполняет дефицит. Её инфинитезимальное
действие имеет вид
\begin{equation}
 \delta_\alpha\Sigma=\rho(\alpha)\Sigma,
\end{equation}
и потому при $\Sigma=0$ имеет нулевой translational tangent. Линейная
representation orbit и аддитивная shift orbit являются разными объектами.

Полную symmetry можно построить условно, добавив новую Stückelberg-копию
$Y_\Sigma$ и используя разность
\begin{equation}
 D(\Sigma,Y_\Sigma)=\Sigma-Y_\Sigma,
 \qquad
 (\Sigma,Y_\Sigma)\mapsto
 (\Sigma+\varepsilon,Y_\Sigma+\varepsilon).
\end{equation}
Тогда
\begin{equation}
 H_{\rm St}=D^TG_YD
 =\begin{pmatrix}G_Y&-G_Y\\-G_Y&G_Y\end{pmatrix},
 \qquad
 \rank H_{\rm St}=6,
 \qquad
 \dim\ker H_{\rm St}=10,
\end{equation}
а diagonal shift имеет rank $8$ и лежит в $\ker H_{\rm St}$. Gauge
$Y_\Sigma=0$ точно возвращает исходный $G_Y$. Условие
$F=QY_\Sigma$ даёт невырожденный Faddeev--Popov operator
\begin{equation}
 M_{\rm FP}=Q,
 \qquad \rank M_{\rm FP}=8,
 \qquad \det M_{\rm FP}=3969,
\end{equation}
то есть именно требуемый odd determinant.

Эта конструкция показывает согласованность всего Mathai--Quillen-пакета,
но не его происхождение. Inherited injection независимой
$Y_\Sigma$-копии имеет rank ноль. Ввести её с diagonal shift означает
добавить ровно тот новый pure-gauge carrier, существование которого должно
было быть доказано.

\begin{theorem}[Условное Stückelberg-замыкание]
На doubled carrier существует rank-$8$ diagonal shift symmetry,
инвариантный parent $D^TG_YD$ и gauge fixing с FP determinant $3969$.
После факторизации исходный $\Sigma$-Hessian восстанавливается точно.
Условный статус равен $6/7$.
\end{theorem}

\begin{nogocriterion}[Финальный no-go particle-wrinkle ветви]
Текущий parent не имеет требуемой rank-$8$ shift symmetry: его квадратичное
ядро имеет dimension $2$, ordinary gauge tangent при $\Sigma=0$ равен
нулю, а независимая Stückelberg-копия не унаследована. Поэтому полный
Mathai--Quillen parent нельзя считать физически выведенным; его статус
остаётся условным.
\end{nogocriterion}

\begin{protocol}\begin{enumerate}
\item inherited $\Sigma$ Hessian rank/nullity: \textbf{$6/2$};
\item required full shift rank: \textbf{$8$};
\item translation-breaking rank: \textbf{$6$};
\item maximum quadratic flat-shift rank: \textbf{$2$};
\item ordinary gauge translational rank at $\Sigma=0$: \textbf{$0$};
\item conditional diagonal Stückelberg shift rank: \textbf{$8$};
\item conditional parent rank/nullity: \textbf{$6/10$};
\item inherited Stückelberg-copy injection rank: \textbf{$0$};
\item conditional/inherited status: \textbf{$6/7$ и $3/7$};
\item следующий гейт: \textbf{финальное заключение Тома X и программа Тома XI}.
\end{enumerate}\end{protocol}

Машинный сертификат сохранён в
\path{s2t_v10_particle_wrinkle_dislocation_callias_equal_charge_twist_m4_h15_r2_mathai_quillen_shift_symmetry_parent_origin_gate_results.json}.